\section{Chapter Preface}
Mature theories eventually become axiomatic. Mechanics, probability, and topology each evolved from useful techniques into formal systems with explicit primitives, assumptions, and theorems.

CLIO reaches that transition point here. Earlier chapters introduced the pieces; this chapter consolidates them into a coherent tuple-level framework with explicit existence, uniqueness, and ambiguity structure.

The resulting trichotomy of inconsistency, identifiability, and residual ambiguity becomes a reusable diagnostic across later chapters.

\section{Derivations and Formal Results}
Define a CLIO system tuple
\[
\mathfrak C=(\stateSpace,\observationSpace,\projection,\constraintSet,F,J).
\]

\begin{theorem}[Existence of preferred admissible reconstruction]
If \(A_y\subseteq\stateSpace\) is nonempty and compact, and \(J:A_y\to\mathbb R\) is continuous, then there exists \(x^\star\in A_y\) such that
\[
J(x^\star)=\min_{x\in A_y}J(x).
\]
\end{theorem}

\begin{proof}
By Weierstrass, a continuous function on a nonempty compact set attains its minimum.
\end{proof}

\begin{proposition}[Uniqueness under strict convexity]
If \(A_y\) is convex and \(J\) is strictly convex on \(A_y\), the minimizer is unique.
\end{proposition}

\begin{theorem}[Admissibility trichotomy]
Exactly one of the following holds:
\[
A_y=\varnothing,\qquad |A_y|=1,\qquad |A_y|>1.
\]
These correspond to inconsistency, identifiability, and residual ambiguity.
\end{theorem}
